% GENERATED FILE. Edit canonical lessons, then run npm run generate:books.
% canonical-source-hash: fnv1a64:ee5aaa3d689ce6a4
% canonical-lessons: UR-W36-dal, UR-W36-rre, UR-W36-do-chashmi-he, UR-W36-gaf

\chapter{Letters With Marks of Their Own}
\label{ch:ur-letters-marks-of-their-own}
\hlchaptermodality{36}

\begin{chapteropening}
\textbf{By the end of this chapter:} \emph{I can write dal, the retroflex re, do-chashmi he and gaf, and find each one in a word I already say.}

Complete the last lesson of chapter 36: i can write dal, the retroflex re, do-chashmi he and gaf, and find each one in a word I already say.
\end{chapteropening}

\section[dāl]{\ur{د} (dāl) — the letter inside \ur{خدا}}

\label{lesson:UR-W36-dal}



\begin{quote}
Before the new one: write \ur{خ} and \ur{چ} once each, and name them.

Now say \textbf{\ur{خدا}} (\emph{khudā}), \textquotedblleft{}God\textquotedblright{}. You have read it for a long time. This lesson takes one letter out of it and puts it in your hand.
\end{quote}

\subsection*{Script you'll notice: \ur{د}}
\textbf{\ur{د}} — \emph{dāl}, the sound \emph{d}.

It is inside \textbf{\ur{خدا}} (\emph{khudā}). Like \textbf{\ur{ا}} and \textbf{\ur{و}}, which you already write, \textbf{\ur{د}} refuses to join leftward: it joins the letter before it, never the one after.

\subsection*{Writing: \ur{د}}
\hlblockfigure{\detokenize{figures/UR-W36-dal-filmstrip.pdf}}{How \ur{د} is written, stroke by stroke}

\begin{itemize}
\raggedright
  \item \textbf{1.} start here: begin at the independent form's upper tip and descend through the folded shoulder
  \item \textbf{2.} turn left along the baseline without lifting — and only now lift
\end{itemize}

\textbf{Pen lifts: 0.} The pen never leaves the paper.

\begin{quote}
Stroke order is one attested teaching order; hands and teachers vary. Source: Zer o Zabar, Dāl, re, and wāw, independent \ur{د} handwriting animation and Dāl instructions (Northwestern University Libraries Open Textbook, 2026).
\end{quote}

\subsection*{Your turn}
\begin{itemize}
\raggedright
  \item \emph{Look:} at \textbf{\ur{خدا}} and point at \ur{د} inside it
  \item \emph{Trace it:} \ur{د} three times, saying \emph{d} as you finish each one
  \item \emph{Write it:} \ur{چ}, then \ur{د}, from memory
  \item \emph{Read it:} \textbf{\ur{خدا}} aloud once more
\end{itemize}

\begin{tcolorbox}[breakable,colback=teal!4,colframe=teal!35!black,title={Before you move on}]
Which letter is this — \ur{د}? (\textbf{dāl}, \emph{d}.) Which word did it come from? (\textbf{\ur{خدا}}, \emph{khudā}, \textquotedblleft{}God\textquotedblright{}.)
\end{tcolorbox}
\section[ṛe]{\ur{ڑ} (ṛe) — the letter inside \ur{پڑھنا}}

\label{lesson:UR-W36-rre}



\begin{quote}
Before the new one: write \ur{چ} and \ur{د} once each, and name them.

Now say \textbf{\ur{پڑھنا}} (\emph{paṛhnā}), \textquotedblleft{}to read\textquotedblright{}. You have read it for a long time. This lesson takes one letter out of it and puts it in your hand.
\end{quote}

\subsection*{Script you'll notice: \ur{ڑ}}
\textbf{\ur{ڑ}} — \emph{ṛe}, an \emph{r} made by curling the tongue back and flicking it forward.

It is inside \textbf{\ur{پڑھنا}} (\emph{paṛhnā}), \textquotedblleft{}to read\textquotedblright{}. It is \textbf{\ur{ر}} with the same \textbf{small curled-back mark} you put on \textbf{\ur{ٹ}}. Urdu uses one mark for one idea: tongue curled back.

\subsection*{Writing: \ur{ڑ}}
\hlblockfigure{\detokenize{figures/UR-W36-rre-filmstrip.pdf}}{How \ur{ڑ} is written, stroke by stroke}

\begin{itemize}
\raggedright
  \item \textbf{1.} start here: draw the independent re-series body downward
  \item \textbf{2.} without lifting, continue curving to the left
  \item \textbf{3.} after one lift, draw the small retroflex mark downward, back upward, and down again to close its loop — and only now lift
\end{itemize}

\textbf{Pen lifts: 1.}

\begin{quote}
Stroke order is one attested teaching order; hands and teachers vary. Source: Zer o Zabar, Fe, qaf, te, dal, and re, independent \ur{ڑ} calligraphic and handwriting animations and Re instructions (Northwestern University Libraries Open Textbook, 2026).
\end{quote}

\subsection*{Your turn}
\begin{itemize}
\raggedright
  \item \emph{Look:} at \textbf{\ur{پڑھنا}} and point at \ur{ڑ} inside it
  \item \emph{Trace it:} \ur{ڑ} three times, saying \emph{ṛ} as you finish each one
  \item \emph{Write it:} \ur{د}, then \ur{ڑ}, from memory
  \item \emph{Read it:} \textbf{\ur{پڑھنا}} aloud once more
\end{itemize}

\begin{tcolorbox}[breakable,colback=teal!4,colframe=teal!35!black,title={Before you move on}]
Which letter is this — \ur{ڑ}? (\textbf{ṛe}, \emph{ṛ}.) Which word did it come from? (\textbf{\ur{پڑھنا}}, \emph{paṛhnā}, \textquotedblleft{}to read\textquotedblright{}.)
\end{tcolorbox}
\section[do-chashmī he]{\ur{ھ} (do-chashmī he) — the two-eyed he inside \ur{ٹھیک}}

\label{lesson:UR-W36-do-chashmi-he}



\begin{quote}
Before the new one: write \ur{د} and \ur{ڑ} once each, and name them.

Now say \textbf{\ur{ٹھیک}} (\emph{ṭhīk}), \textquotedblleft{}fine, well\textquotedblright{}. You have read it for a long time. This lesson takes one letter out of it and puts it in your hand.
\end{quote}

\subsection*{Script you'll notice: \ur{ھ}}
\textbf{\ur{ھ}} — \emph{do-chashmī he}, \textquotedblleft{}two-eyed \emph{he}\textquotedblright{}.

It is inside \textbf{\ur{ٹھیک}} (\emph{ṭhīk}) and \textbf{\ur{پڑھنا}} (\emph{paṛhnā}). It almost never makes a sound of its own. It rides on the letter before it and adds a puff of breath: \textbf{\ur{ٹ}} + \textbf{\ur{ھ}} is \emph{ṭh}, \textbf{\ur{ڑ}} + \textbf{\ur{ھ}} is \emph{ṛh}. Two round loops, the \textquotedblleft{}two eyes\textquotedblright{}.

\subsection*{Writing: \ur{ھ}}
\hlblockfigure{\detokenize{figures/UR-W36-do-chashmi-he-filmstrip.pdf}}{How \ur{ھ} is written, stroke by stroke}

\begin{itemize}
\raggedright
  \item \textbf{1.} start here: circle the right eye clockwise from the upper center
  \item \textbf{2.} continue down and left along the baseline without lifting
  \item \textbf{3.} without lifting, reverse at the left edge and rise around the left eye
  \item \textbf{4.} without lifting, close at the center and finish with the low leftward sweep — and only now lift
\end{itemize}

\textbf{Pen lifts: 0.} The pen never leaves the paper.

\begin{quote}
Stroke order is one attested teaching order; hands and teachers vary. Source: Zer o Zabar, Chhoṭī he, do-chashmī he, chhoṭī ye, baṛī ye, and punctuation, independent \ur{ھ} calligraphic and handwriting animations and Aspiration with do-chashmī he instructions (Northwestern University Libraries Open Textbook, 2026).
\end{quote}

\subsection*{Your turn}
\begin{itemize}
\raggedright
  \item \emph{Look:} at \textbf{\ur{ٹھیک}} and point at \ur{ھ} inside it
  \item \emph{Trace it:} \ur{ھ} three times, saying \emph{h} as you finish each one
  \item \emph{Write it:} \ur{ڑ}, then \ur{ھ}, from memory
  \item \emph{Read it:} \textbf{\ur{ٹھیک}} aloud once more
\end{itemize}

\begin{tcolorbox}[breakable,colback=teal!4,colframe=teal!35!black,title={Before you move on}]
Which letter is this — \ur{ھ}? (\textbf{do-chashmī he}, \emph{h}.) Which word did it come from? (\textbf{\ur{ٹھیک}}, \emph{ṭhīk}, \textquotedblleft{}fine, well\textquotedblright{}.)
\end{tcolorbox}
\section[gāf]{\ur{گ} (gāf) — the letter that opens \ur{گرمی}}

\label{lesson:UR-W36-gaf}



\begin{quote}
Before the new one: write \ur{ڑ} and \ur{ھ} once each, and name them.

Now say \textbf{\ur{گرمی}} (\emph{garmī}), \textquotedblleft{}heat\textquotedblright{}. You have read it for a long time. This lesson takes one letter out of it and puts it in your hand.
\end{quote}

\subsection*{Script you'll notice: \ur{گ}}
\textbf{\ur{گ}} — \emph{gāf}, the sound \emph{g}.

It opens \textbf{\ur{گرمی}} (\emph{garmī}), \textquotedblleft{}heat\textquotedblright{}. It is \textbf{\ur{ک}}, which you already write, with a \textbf{second slash} above the first.

\subsection*{Writing: \ur{گ}}
\hlblockfigure{\detokenize{figures/UR-W36-gaf-filmstrip.pdf}}{How \ur{گ} is written, stroke by stroke}

\begin{itemize}
\raggedright
  \item \textbf{1.} start here: draw the independent stem downward
  \item \textbf{2.} flow right to left through the flatter bowl and finish with the hook without lifting
  \item \textbf{3.} after one lift, draw the long slash down from the upper right toward the stem
  \item \textbf{4.} after another lift, draw the shorter floating slash above the first — and only now lift
\end{itemize}

\textbf{Pen lifts: 2.}

\begin{quote}
Stroke order is one attested teaching order; hands and teachers vary. Source: Zer o Zabar, Pe, gāf, alif, and lām, independent \ur{گ} calligraphic and handwriting animations and Gāf instructions (Northwestern University Libraries Open Textbook, 2026).
\end{quote}

\subsection*{Your turn}
\begin{itemize}
\raggedright
  \item \emph{Look:} at \textbf{\ur{گرمی}} and point at \ur{گ} inside it
  \item \emph{Trace it:} \ur{گ} three times, saying \emph{g} as you finish each one
  \item \emph{Write it:} \ur{ھ}, then \ur{گ}, from memory
  \item \emph{Read it:} \textbf{\ur{گرمی}} aloud once more
\end{itemize}

\begin{tcolorbox}[breakable,colback=teal!4,colframe=teal!35!black,title={Before you move on}]
Which letter is this — \ur{گ}? (\textbf{gāf}, \emph{g}.) Which word did it come from? (\textbf{\ur{گرمی}}, \emph{garmī}, \textquotedblleft{}heat\textquotedblright{}.)
\end{tcolorbox}
